\begin{definition} \rm 
Let $X$ be a metric space which consists of $n$ points and let $\ell_1,\dots, \ell_N$ be the edge lengths, where $N={n\choose2}$. 
A map 
\[A\colon \left\{\{i,j\}\,|\,1\le i<j\le n \right\}\to\{\ell_1,\dots,\ell_N\}\]
is called {\em combinatorial data} (of edge lengths) when there exists a labeling of points in $X$ such that $A(\{i,j\})=d(P_i,P_j)$ for any $i,j$ $(i<j)$. 

We write $d_{i,j}=d(P_i,P_j)$ in what follows. 
\end{definition}

A {\em multiset} is a set with multiplicity. We will use the symbol $[\>]$ for multisets. For example, although $\{a,a,b\}$ and $\{a,b\}$ are same as a set, $[a,a,b]$ and $[a,b]$ are different multisets. 

\begin{definition}\label{def_2} \rm 
We say that a finite metric space $X$ is {\em determined by the magnitude function} if the multiset of the edge lengths and the combinatorial data are obtained from the magnitude function $M_X(t)$. 
%
%
%
\end{definition}

Recall that the cardinality of a finite metric space is obtained from the magnitude function by $\#X=\lim_{t\to\infty}M_X(t)$ (\cite{L13} Proposition 2.2.6, \cite{LW} Theorem 3). 

\begin{definition} \rm 
A finite metric space is, respectively, {\em rationally independent} (ri); {\em $p$-generic} ($\mbox{g}_p$) for $p$ a natural number; or satisfying the {\em strict virtual triangle inequality} (svti) if the following condition is satisfied respectively:
\begin{enumerate}
\item[($\mbox{ri}$)] The edge lengths are linearly independent over $\QQ$. In other words, the sums of edge lengths do not match for different combinations with multiplicity.
%
\item[($\mbox{g}_p$)] The sums of edge lengths do not match for different combinations of $p$ or fewer edges with multiplicity.
%
\item[(svti)] $\displaystyle \max_{1\le i,j\le n} d_{ij}<2\min_{1\le k,l\le n, k\ne l} d_{kl}$ \qquad ($n=\#X$).
\end{enumerate}
\end{definition}

Note that the rational independence implies $p$-genericity for any $p$. 
We remark that our conditions are not well suited for graphs. Any graph with graph metric with two edges or more is rationally dependent. Any connected graph except for complete graphs does not satisfy the strict virtual triangle inequality condition. 
A set consisting of $n$-simplex vertices close to a regular $n$-simplex in $\RR^{n+1}$ satisfies the strict virtual triangle inequality condition, whereas a set consisting of the vertices of a needle-shaped tetrahedron does not satisfy this condition. 


\begin{theorem}\label{theorem}
A finite metric space $X$ 
is determined by the magnitude function if $X$ satisfies one of the following conditions ($n=\#X$). 

\begin{enumerate}
\item $n=3$. 
\item $X$ is rationally independent. 
\item $X$ is $5$-generic and satisfies the strict virtual triangle inequality condition. %
\item $n=4$ and $X$ satisfies the strict virtual triangle inequality condition. %
\end{enumerate}

\end{theorem}
